\subsubsection{La-MAML, C-MAML, and S-MAML}

\gls{lamaml} \cite{ll_lamaml_2020} is a meta-learning-based \gls{cl} method that mitigates \gls{cf} by maintaining a set of per-parameter adaptive learning rates that are meta-trained to minimise forgetting on a replay buffer of past data. Two ablation variants, \gls{cmaml} and \gls{smaml}, share the same inner-loop structure but differ in whether the learning rates are trained and how the outer weight update is applied.

\paragraph{Background: Model-Agnostic Meta-Learning (MAML)}

\gls{maml} \cite{maml_2017} frames learning as a bi-level optimisation problem. The \emph{inner loop} simulates rapid task adaptation: starting from the current model parameters $\theta$, a small number of gradient steps on a task-specific loss produce \emph{fast weights} $\theta'$ that are specialised for that task. The \emph{outer loop} then evaluates a meta-objective on held-out data using $\theta'$ and back-propagates its gradient all the way through the inner update steps to adjust $\theta$. Because the meta-gradient accounts for how the inner update will change the parameters, optimising $\theta$ in this way finds an initialisation from which fast adaptation is possible, rather than one that simply performs well before any adaptation.

In the standard MAML formulation the inner step uses a fixed scalar learning rate $\alpha$ and the outer step uses a separate optimiser with rate $\beta$. The two rates play distinct roles: $\alpha$ controls how far the fast weights move from the initialisation, while $\beta$ controls how the shared initialisation itself evolves in response to the meta-gradient. La-MAML inherits this two-level structure and applies it to \gls{cl} by replacing the held-out task data used for the meta-objective with a replay buffer of past observations, turning the outer loop into an explicit anti-forgetting signal.

\paragraph{Per-parameter adaptive learning rates}

All three variants augment the shared backbone with a per-parameter scalar $\alpha_i \geq 0$ for each network parameter $\theta_i$, initialised to a small constant $\alpha_\text{init}$. In all variants $\alpha_i$ serves as the inner-loop step size; whether it is held fixed or meta-trained, and whether it additionally governs the outer weight update, differs across variants and is described in the per-variant sections below.

\paragraph{Replay buffer}

All three variants maintain a fixed-capacity replay buffer $\mathcal{M}$ of size $K$ that is updated via reservoir sampling (identical to \gls{er-res}; see \eqref{eqn:reservoir}). A mixed batch $B^\text{rep} = (x^\text{rep}, y^\text{rep})$ drawn from $\mathcal{M}$ is used as the \emph{outer} (meta) objective to assess generalisation after each inner update.

\paragraph{Inner update}

At each training step the current mini-batch of size $B$ is divided into $K_\text{meta}$ equal sub-batches. For sub-batch $i$ the model parameters are updated \emph{in-place in a fast-weight copy} using a single gradient step scaled by the per-parameter rates:
%
\begin{equation}\label{eqn:lamaml_inner}
    \theta'_j
    =
    \theta_j
    -
    \alpha_j \cdot
    \frac{\partial \mathcal{L}_\text{CE}(B_i;\,\theta)}
         {\partial \theta_j},
    \qquad
    \alpha_j = \mathrm{ReLU}(\hat\alpha_j),
\end{equation}
%
where the ReLU enforces non-negative step sizes and $\hat\alpha_j$ is the raw (possibly negative) stored parameter. The computational graph through the inner step is retained so that the meta-loss can differentiate through it via second-order gradients.

\paragraph{Meta-objective (outer loss)}

After each inner step, the fast weights $\theta'$ are evaluated on the replay batch to measure how well the proposed update generalises across all seen tasks:
%
\begin{equation}\label{eqn:lamaml_meta_loss}
    \mathcal{L}_\text{meta}
    =
    \frac{1}{K_\text{meta}}
    \sum_{i=1}^{K_\text{meta}}
    \mathcal{L}_\text{CE}\!\left(B^\text{rep};\, \theta'_i\right).
\end{equation}
%
The gradient $\partial \mathcal{L}_\text{meta} / \partial \hat\alpha_j$ propagates back through the inner update \eqref{eqn:lamaml_inner} to shape the per-parameter step sizes.

\paragraph{La-MAML: asynchronous weight update}

\gls{lamaml} performs the outer update in two sequential steps. First, the per-parameter learning rates are updated via SGD with step size $\eta_\alpha$:
%
\begin{equation}\label{eqn:lamaml_lr_update}
    \hat\alpha_j
    \leftarrow
    \hat\alpha_j
    -
    \eta_\alpha \cdot
    \frac{\partial \mathcal{L}_\text{meta}}{\partial \hat\alpha_j}.
\end{equation}
%
The network weights are then updated \emph{using the just-updated learning rates as per-parameter step sizes} applied to the gradient of $\mathcal{L}_\text{meta}$ with respect to the current (non-fast) weights:
%
\begin{equation}\label{eqn:lamaml_weight_update}
    \theta_j
    \leftarrow
    \theta_j
    -
    \mathrm{ReLU}(\hat\alpha_j) \cdot
    \frac{\partial \mathcal{L}_\text{meta}}{\partial \theta_j}.
\end{equation}
%
This coupling means the backbone immediately benefits from the freshly meta-updated step sizes, rather than a fixed global learning rate.

\paragraph{\gls{cmaml}}

\gls{cmaml} is the closest variant to vanilla MAML \cite{} applied in the \gls{cl} setting. The per-parameter scalars $\hat\alpha_j$ are fixed at $\alpha_\text{init}$ throughout training. They serve exclusively as the inner-loop step size in \eqref{eqn:lamaml_inner}: they control how far the fast weights move and, through the retained second-order graph, shape the meta-gradient, but they play no role in the outer update. The outer weight update uses a standard SGD optimiser with momentum and a single global step size $\beta$:
%
\begin{equation}\label{eqn:cmaml_weight_update}
    \theta_j
    \leftarrow
    \theta_j
    -
    \beta \cdot
    \frac{\partial \mathcal{L}_\text{meta}}{\partial \theta_j}.
\end{equation}
%
\gls{cmaml} therefore has two separate, non-interacting hyperparameters: $\alpha_\text{init}$ governs inner adaptation depth and $\beta$ governs outer plasticity. Neither is informed by the other.

\paragraph{\gls{smaml}: learned inner rates, fixed outer rate}

\gls{smaml} adds meta-learning of the per-parameter rates on top of \gls{cmaml}, but the outer weight update still uses a fixed global step size $\beta$ as in \eqref{eqn:cmaml_weight_update}. The learning rates are updated by \eqref{eqn:lamaml_lr_update} and then the weights are updated with $\beta$. The freshly-updated $\hat\alpha_j$ values are not consulted again for the weight step; they only take effect in the \emph{next} training step's inner update.

This means \gls{smaml} uses $\hat\alpha_j$ in only one place (the inner update), whereas \gls{lamaml} uses the same $\hat\alpha_j$ in both the inner update and the outer weight step \eqref{eqn:lamaml_weight_update}. Comparing \gls{smaml} to \gls{lamaml} therefore isolates exactly the effect of feeding the learned rates back into the outer update. Note that, as an implementation artefact, the SGD optimiser used for outer weight updates in \gls{cmaml} and \gls{smaml} carries momentum ($0.9$) whereas the manual per-parameter update in \gls{lamaml} does not; this partially confounds the comparison.

\paragraph{Summary of variant differences}

\begin{center}
\begin{tabular}{lcccc}
\hline
Variant & Per-param.\ inner LR & Inner LR meta-trained & Outer weight step size & Outer step uses learned LR \\
\hline
C-MAML  & Yes (fixed $\alpha_\text{init}$) & No  & global scalar $\beta$ & No \\
S-MAML  & Yes (learned $\alpha_j$)         & Yes & global scalar $\beta$ & No \\
La-MAML & Yes (learned $\alpha_j$)         & Yes & per-param.\ $\alpha_j$ & Yes \\
\hline
\end{tabular}
\end{center}

\paragraph{Inference}

All three variants use a standard linear classifier at inference time. In \gls{til}, task identity is used to select the relevant task logit slice. In \gls{cil}, prediction is made from the full logit vector across all seen classes.
